\chapter{Introduction}
\label{ch:intro}
\vspace{2em}

% TODO: One or two opening paragraphs setting the scene before the overview.

\section{Overview}
\label{intro:sec:overview}

% TODO: Overview of the thesis -- the virtual teaching assistant, its
% deployment context (NUS Canvas course materials), and the structure of the
% document.

\section{Problem Statement}
\label{intro:sec:problem}

Educators operate under a fixed time budget divided across teaching,
preparation, marking, administration, and --- in higher education ---
research. Survey evidence shows that non-teaching duties consume a large and
stable share of this budget: Singapore teachers report four hours per week of
administrative work (above the OECD average of three) within a 47.3-hour
working week of which only 17.7 hours is teaching, and 53\% rank excessive
administrative work as their foremost source of occupational
stress~\cite{oecd25:TalisSG}. Panel evidence from higher education shows that
this burden is a quality problem, not merely a workload problem: non-teaching
duties that carry no compensating load relief measurably reduce teaching
quality, while research time feeds back positively into teaching quality over
most of its observed range~\cite{apec15:GarciaGallego}. The AI-in-education
literature has long proposed intelligent systems --- chatbots, automated
assessment, learning analytics --- as the remedy, arguing that such
applications \enquote{minimize the administrative tasks of a teacher to
invest more in teaching and guiding students}~\cite{sust22:Ahmad}; yet that
same literature concedes that its claims remain qualitative and
quantitatively unverified, and the systems it surveys predate modern
retrieval-augmented large language models~\cite{sust22:Ahmad,
ijethe19:ZawackiRichter}.

This thesis presents a retrieval-augmented virtual teaching assistant built
over NUS Canvas course materials that absorbs routine student-support load
--- clarification, content navigation, and logistics queries --- and, unlike
prior work, instruments the interaction telemetry required to quantify the
instructor-hours it returns, connecting the workload-measurement,
teaching-quality, and AI-intervention literatures for the first time
(\autoref{ch:review}).

\section{Objectives}
\label{intro:sec:objectives}

% TODO: Concrete objectives -- e.g. (i) build a RAG-based teaching assistant
% over institutional course materials, (ii) make it deployable on modest
% hardware (context compression), (iii) instrument query deflection to
% estimate instructor-hours returned.

\section{Summary of Contributions}
\label{intro:sec:contributions}

% TODO: Summary of contributions, mirroring the gap statement in
% \autoref{review:subsec:gap}.
